
\lecture
{17}
{Wednesday, 7 October 2026}
{Introduction to Probability}
{Dr. Nishant Chandgotia}


% ============================================================
\section{Subsequence Characterization of Convergence in Probability}
% ============================================================

\begin{theorem}
{Subsequence Characterization of Convergence in Probability}
{subsequence-characterization-convergence-probability}

Let $X,X_1,X_2,\ldots$ be random variables. Then
\[
    X_n \xrightarrow{P} X
\]
if and only if every subsequence $(X_{n_k})$ has a further
subsequence $(X_{n_{k_j}})$ such that
\[
    X_{n_{k_j}} \xrightarrow{\mathrm{a.s.}} X.
\]

\end{theorem}


\begin{proof}

We prove the two implications separately.

\medskip

\noindent
\textbf{($\Rightarrow$) Suppose that $X_n\xrightarrow{P}X$.}

We have to prove that every subsequence of $(X_n)$ has a further
subsequence which converges to $X$ almost surely.

Let
\[
    (X_{n_k})
\]
be an arbitrary subsequence of $(X_n)$.

Since
\[
    X_n\xrightarrow{P}X,
\]
by definition, for every $\varepsilon>0$,
\[
    \mathbb P(|X_n-X|>\varepsilon)\longrightarrow0.
\]

We shall construct a further subsequence
\[
    X_{n_{k_1}},X_{n_{k_2}},\ldots
\]
such that
\[
    \mathbb P\left(
        |X_{n_{k_j}}-X|>\frac{1}{2^j}
    \right)
    <\frac{1}{2^j}
\]
for every $j\geq1$.

To see that this is possible, first choose $k_1$ such that
\[
    \mathbb P\left(
        |X_{n_{k_1}}-X|>1
    \right)<\frac12.
\]
This is possible because
\[
    \mathbb P(|X_n-X|>1)\longrightarrow0.
\]

Having chosen $k_1,\ldots,k_{j-1}$, choose $k_j>k_{j-1}$ such that
\[
    \mathbb P\left(
        |X_{n_{k_j}}-X|>\frac{1}{2^j}
    \right)
    <\frac{1}{2^j}.
\]
Again, this is possible because
\[
    \mathbb P\left(
        |X_n-X|>\frac{1}{2^j}
    \right)
    \longrightarrow0.
\]

Thus we have constructed a further subsequence satisfying
\[
    \mathbb P\left(
        |X_{n_{k_j}}-X|>\frac{1}{2^j}
    \right)
    <\frac{1}{2^j}.
\]

Define the events
\[
    A_j
    =
    \left\{
        |X_{n_{k_j}}-X|>\frac{1}{2^j}
    \right\}.
\]

Then
\[
    \sum_{j=1}^{\infty}\mathbb P(A_j)
    \leq
    \sum_{j=1}^{\infty}\frac{1}{2^j}
    =1<\infty.
\]

Therefore, by the first Borel--Cantelli lemma,
\[
    \mathbb P(A_j\text{ i.o.})=0.
\]

Hence, with probability one, only finitely many of the events $A_j$
occur.

Thus, for almost every $\omega$, there exists $J=J(\omega)$ such that
for all $j\geq J$,
\[
    |X_{n_{k_j}}(\omega)-X(\omega)|
    \leq\frac{1}{2^j}.
\]

Since
\[
    \frac{1}{2^j}\longrightarrow0,
\]
we obtain
\[
    X_{n_{k_j}}(\omega)\longrightarrow X(\omega)
\]
for almost every $\omega$.

Therefore,
\[
    \boxed{
    X_{n_{k_j}}\xrightarrow{\mathrm{a.s.}}X.
    }
\]

Since the original subsequence $(X_{n_k})$ was arbitrary, every
subsequence has a further subsequence which converges almost surely
to $X$.

\medskip

\noindent
\textbf{($\Leftarrow$) Suppose that every subsequence of $(X_n)$ has
a further subsequence which converges almost surely to $X$.}

We prove that
\[
    X_n\xrightarrow{P}X.
\]

Suppose, to the contrary, that
\[
    X_n\not\xrightarrow{P}X.
\]

By the definition of convergence in probability, this means that there
exists some $\varepsilon_0>0$ for which
\[
    \mathbb P(|X_n-X|>\varepsilon_0)
\]
does not converge to $0$.

Therefore, there exists some $\delta>0$ and a subsequence
\[
    X_{n_k}
\]
such that
\[
    \mathbb P(|X_{n_k}-X|>\varepsilon_0)
    \geq\delta
\]
for every $k\geq1$.

Indeed, if no such $\delta$ and subsequence existed, then
\[
    \mathbb P(|X_n-X|>\varepsilon_0)\longrightarrow0,
\]
which would contradict our assumption.

By the hypothesis of the theorem, this subsequence $(X_{n_k})$ has a
further subsequence
\[
    X_{n_{k_j}}
\]
such that
\[
    X_{n_{k_j}}\xrightarrow{\mathrm{a.s.}}X.
\]

We now show that almost sure convergence implies convergence in
probability.

Since
\[
    X_{n_{k_j}}\xrightarrow{\mathrm{a.s.}}X,
\]
for almost every $\omega$,
\[
    |X_{n_{k_j}}(\omega)-X(\omega)|
    \longrightarrow0.
\]

Fix $\varepsilon_0>0$ and define
\[
    B_m
    =
    \bigcup_{j=m}^{\infty}
    \left\{
        |X_{n_{k_j}}-X|>\varepsilon_0
    \right\}.
\]

The events satisfy
\[
    B_1\supseteq B_2\supseteq B_3\supseteq\cdots.
\]

Moreover, since
\[
    X_{n_{k_j}}\xrightarrow{\mathrm{a.s.}}X,
\]
only finitely many of the events
\[
    \left\{
        |X_{n_{k_j}}-X|>\varepsilon_0
    \right\}
\]
occur almost surely. Hence
\[
    \mathbb P\left(\bigcap_{m=1}^{\infty}B_m\right)=0.
\]

By continuity of probability from above,
\[
    \mathbb P(B_m)
    \longrightarrow
    \mathbb P\left(\bigcap_{m=1}^{\infty}B_m\right)
    =0.
\]

But for every $j\geq m$,
\[
    \left\{
        |X_{n_{k_j}}-X|>\varepsilon_0
    \right\}
    \subseteq B_m.
\]

Therefore,
\[
    \mathbb P(|X_{n_{k_j}}-X|>\varepsilon_0)
    \leq \mathbb P(B_m).
\]

For fixed $m$, this holds for every $j\geq m$. Hence
\[
    \limsup_{j\to\infty}
    \mathbb P(|X_{n_{k_j}}-X|>\varepsilon_0)
    \leq \mathbb P(B_m).
\]

Letting $m\to\infty$, we get
\[
    \limsup_{j\to\infty}
    \mathbb P(|X_{n_{k_j}}-X|>\varepsilon_0)
    =0.
\]

Thus
\[
    \mathbb P(|X_{n_{k_j}}-X|>\varepsilon_0)
    \longrightarrow0.
\]

However, $(X_{n_{k_j}})$ is a further subsequence of $(X_{n_k})$, and
we constructed $(X_{n_k})$ so that
\[
    \mathbb P(|X_{n_k}-X|>\varepsilon_0)
    \geq\delta
\]
for every $k$.

Consequently,
\[
    \mathbb P(|X_{n_{k_j}}-X|>\varepsilon_0)
    \geq\delta
\]
for every $j$,
which contradicts
\[
    \mathbb P(|X_{n_{k_j}}-X|>\varepsilon_0)
    \longrightarrow0.
\]

This contradiction shows that our assumption
\[
    X_n\not\xrightarrow{P}X
\]
was false. Therefore,
\[
    \boxed{
    X_n\xrightarrow{P}X.
    }
\]

Hence,
\[
    \boxed{
    X_n\xrightarrow{P}X
    \quad\Longleftrightarrow\quad
    \text{every subsequence of $(X_n)$ has a further subsequence
    converging a.s. to $X$.}
    }
\]

\end{proof}


% ============================================================
\section{Infinite First Moment}
% ============================================================

\begin{theorem}
{Durrett's Theorem 2.3.8}
{durrett-theorem-2-3-8}

Let $X_1,X_2,\ldots$ be independent and identically distributed random
variables such that
\[
    \mathbb{E}|X_i|=\infty.
\]

Then
\[
    \mathbb{P}\bigl(|X_n|\geq n \text{ i.o.}\bigr)=1.
\]

Consequently, if
\[
    S_n=X_1+\cdots+X_n,
\]
then
\[
    \mathbb{P}\left(
        \lim_{n\to\infty}\frac{S_n}{n}
        \text{ exists in }(-\infty,\infty)
    \right)=0.
\]

\end{theorem}


\begin{proof}

We divide the proof into two parts.

\medskip

\noindent
\textbf{Part I: Proving that $|X_n|\geq n$ infinitely often almost surely.}

For any nonnegative random variable $Y$, the tail-integral formula gives
\[
    \mathbb{E}Y
    =
    \int_0^\infty \mathbb{P}(Y>x)\,dx.
\]

Applying this with $Y=|X_1|$, we obtain
\[
    \mathbb{E}|X_1|
    =
    \int_0^\infty
    \mathbb{P}(|X_1|>x)\,dx.
\]

Now divide the integral into intervals $[n,n+1)$:
\[
\begin{aligned}
    \mathbb{E}|X_1|
    &=
    \sum_{n=0}^\infty
    \int_n^{n+1}
    \mathbb{P}(|X_1|>x)\,dx.
\end{aligned}
\]

Since the function
\[
    x\longmapsto\mathbb{P}(|X_1|>x)
\]
is nonincreasing, for $x\in[n,n+1)$ we have
\[
    \mathbb{P}(|X_1|>x)
    \leq
    \mathbb{P}(|X_1|>n).
\]

Therefore,
\[
\begin{aligned}
    \mathbb{E}|X_1|
    &\leq
    \sum_{n=0}^\infty
    \int_n^{n+1}
    \mathbb{P}(|X_1|>n)\,dx
    \\
    &=
    \sum_{n=0}^\infty
    \mathbb{P}(|X_1|>n).
\end{aligned}
\]

But by assumption,
\[
    \mathbb{E}|X_1|=\infty.
\]

Hence
\[
    \sum_{n=0}^\infty
    \mathbb{P}(|X_1|>n)
    =\infty.
\]

Since the random variables $X_1,X_2,\ldots$ are identically distributed,
\[
    \sum_{n=1}^\infty
    \mathbb{P}(|X_n|>n)
    =
    \sum_{n=1}^\infty
    \mathbb{P}(|X_1|>n)
    =\infty.
\]

Define the events
\[
    A_n=\{|X_n|>n\}.
\]

Because $X_1,X_2,\ldots$ are independent, the events
$A_1,A_2,\ldots$ are independent.

Therefore, by the second Borel--Cantelli lemma,
\[
    \sum_{n=1}^\infty\mathbb{P}(A_n)=\infty
    \quad\Longrightarrow\quad
    \mathbb{P}(A_n\text{ i.o.})=1.
\]

Consequently,
\[
    \mathbb{P}\bigl(|X_n|>n\text{ i.o.}\bigr)=1.
\]

Since
\[
    \{|X_n|>n\text{ i.o.}\}
    \subseteq
    \{|X_n|\geq n\text{ i.o.}\},
\]
we obtain
\[
    \boxed{
    \mathbb{P}\bigl(|X_n|\geq n\text{ i.o.}\bigr)=1.
    }
    \]


\medskip

\noindent
\textbf{Part II: Proving that $S_n/n$ cannot converge to a finite limit.}

Let
\[
    C
    =
    \left\{
        \omega:
        \lim_{n\to\infty}\frac{S_n(\omega)}{n}
        \text{ exists in }(-\infty,\infty)
    \right\}.
\]

We will prove that
\[
    \mathbb{P}(C)=0.
\]

Suppose, for contradiction, that
\[
    \omega\in C.
\]

Then, by definition, there exists some finite real number $L=L(\omega)$
such that
\[
    \frac{S_n(\omega)}{n}\longrightarrow L.
\]

In particular, the sequence $S_n(\omega)/n$ is bounded. Hence there
exists a finite constant $M=M(\omega)$ such that
\[
    \left|\frac{S_n(\omega)}{n}\right|\leq M
\]
for all sufficiently large $n$.

We now consider the difference between two consecutive terms:
\[
    \frac{S_n}{n}
    -
    \frac{S_{n+1}}{n+1}.
\]

Since
\[
    S_{n+1}=S_n+X_{n+1},
\]
we have
\[
\begin{aligned}
    \frac{S_n}{n}
    -
    \frac{S_{n+1}}{n+1}
    &=
    \frac{S_n}{n}
    -
    \frac{S_n+X_{n+1}}{n+1}
    \\
    &=
    \frac{(n+1)S_n-n(S_n+X_{n+1})}
         {n(n+1)}
    \\
    &=
    \frac{S_n-nX_{n+1}}
         {n(n+1)}
    \\
    &=
    \frac{S_n}{n(n+1)}
    -
    \frac{X_{n+1}}{n+1}.
\end{aligned}
\]

Thus
\[
    \boxed{
    \frac{S_n}{n}
    -
    \frac{S_{n+1}}{n+1}
    =
    \frac{S_n}{n(n+1)}
    -
    \frac{X_{n+1}}{n+1}.
    }
\]

Because $\omega\in C$,
\[
    \frac{S_n(\omega)}{n}\longrightarrow L.
\]

Therefore,
\[
    \frac{S_n(\omega)}{n(n+1)}
    =
    \frac{1}{n+1}
    \frac{S_n(\omega)}{n}
    \longrightarrow0.
\]

Hence
\[
    \boxed{
    \frac{S_n(\omega)}{n(n+1)}
    \longrightarrow0.
    }
\]

Now define
\[
    A
    =
    \left\{
        \omega:
        |X_n(\omega)|\geq n
        \text{ infinitely often}
    \right\}.
\]

From Part I,
\[
    \mathbb{P}(A)=1.
\]

Suppose now that
\[
    \omega\in A\cap C.
\]

Since $\omega\in A$,
\[
    |X_n(\omega)|\geq n
\]
for infinitely many $n$.

Replacing $n$ by $n+1$, this means
\[
    |X_{n+1}(\omega)|\geq n+1
\]
for infinitely many $n$.

Consequently, for infinitely many $n$,
\[
    \frac{|X_{n+1}(\omega)|}{n+1}\geq1.
\]

Using the identity above and the reverse triangle inequality,
\[
\begin{aligned}
\left|
    \frac{S_n(\omega)}{n}
    -
    \frac{S_{n+1}(\omega)}{n+1}
\right|
&=
\left|
    \frac{S_n(\omega)}{n(n+1)}
    -
    \frac{X_{n+1}(\omega)}{n+1}
\right|
\\
&\geq
\frac{|X_{n+1}(\omega)|}{n+1}
-
\frac{|S_n(\omega)|}{n(n+1)}.
\end{aligned}
\]

For infinitely many $n$,
\[
    \frac{|X_{n+1}(\omega)|}{n+1}\geq1.
\]

Moreover,
\[
    \frac{|S_n(\omega)|}{n(n+1)}
    \longrightarrow0.
\]

Therefore, for all sufficiently large $n$ among those infinitely many
indices,
\[
    \frac{|S_n(\omega)|}{n(n+1)}
    <\frac13.
\]

Hence, infinitely often,
\[
\begin{aligned}
\left|
    \frac{S_n(\omega)}{n}
    -
    \frac{S_{n+1}(\omega)}{n+1}
\right|
&\geq
1-\frac13
\\
&=\frac23.
\end{aligned}
\]

Thus
\[
    \boxed{
    \left|
    \frac{S_n(\omega)}{n}
    -
    \frac{S_{n+1}(\omega)}{n+1}
    \right|
    >\frac23
    \quad\text{i.o.}
    }
\]

However, this contradicts $\omega\in C$.

Indeed, if
\[
    \frac{S_n(\omega)}{n}\longrightarrow L,
\]
then also
\[
    \frac{S_{n+1}(\omega)}{n+1}\longrightarrow L.
\]

Therefore,
\[
\begin{aligned}
\left|
    \frac{S_n(\omega)}{n}
    -
    \frac{S_{n+1}(\omega)}{n+1}
\right|
&\leq
\left|
    \frac{S_n(\omega)}{n}-L
\right|
+
\left|
    \frac{S_{n+1}(\omega)}{n+1}-L
\right|
\\
&\longrightarrow0.
\end{aligned}
\]

Thus the difference must converge to $0$, which is impossible if it is
larger than $2/3$ infinitely often.

We have therefore proved that
\[
    A\cap C=\varnothing.
\]

Since
\[
    \mathbb{P}(A)=1,
\]
we have
\[
    \mathbb{P}(A^c)=0.
\]

Because
\[
    C\subseteq A^c,
\]
it follows that
\[
    \mathbb{P}(C)
    \leq
    \mathbb{P}(A^c)
    =0.
\]

Hence
\[
    \boxed{
    \mathbb{P}\left(
        \lim_{n\to\infty}\frac{S_n}{n}
        \text{ exists in }(-\infty,\infty)
    \right)=0.
    }
\]

This completes the proof.

\end{proof}


% ============================================================
\section{Strong Law for Pairwise Independent Events}
% ============================================================

\begin{theorem}
{Strong Law for Pairwise Independent Events}
{strong-law-pairwise-independent-events}

If $A_1,A_2,\ldots$ are pairwise independent and
\[
    \sum_{n=1}^{\infty}\mathbb P(A_n)=\infty,
\]
then
\[
    \boxed{
    \frac{\displaystyle\sum_{m=1}^{n}\mathbf{1}_{A_m}}
         {\displaystyle\sum_{m=1}^{n}\mathbb P(A_m)}
    \xrightarrow{\mathrm{a.s.}}1.
    }
\]

\end{theorem}


\begin{proof}

Let
\[
    X_m=\mathbf{1}_{A_m},
\]
and define
\[
    S_n=X_1+\cdots+X_n
    =
    \sum_{m=1}^{n}\mathbf{1}_{A_m}.
\]

Since the events $A_m$ are pairwise independent, the random variables
$X_m=\mathbf{1}_{A_m}$ are pairwise independent. Hence they are
uncorrelated, and therefore
\[
    \operatorname{Var}(S_n)
    =
    \sum_{m=1}^{n}\operatorname{Var}(X_m).
\]

Now $X_m\in\{0,1\}$, so
\[
    X_m^2=X_m.
\]

Therefore,
\[
    \mathbb E(X_m^2)=\mathbb E(X_m)
\]
and hence
\[
\begin{aligned}
    \operatorname{Var}(X_m)
    &=
    \mathbb E(X_m^2)
    -
    \bigl(\mathbb E X_m\bigr)^2
    \\
    &=
    \mathbb E(X_m)
    -
    \bigl(\mathbb E X_m\bigr)^2
    \\
    &\leq
    \mathbb E(X_m).
\end{aligned}
\]

Thus
\[
\begin{aligned}
    \operatorname{Var}(S_n)
    &=
    \sum_{m=1}^{n}\operatorname{Var}(X_m)
    \\
    &\leq
    \sum_{m=1}^{n}\mathbb E(X_m)
    \\
    &=
    \mathbb E(S_n).
\end{aligned}
\]

Since
\[
    \mathbb E(X_m)
    =
    \mathbb P(A_m),
\]
we have
\[
    \boxed{
    \mathbb E(S_n)
    =
    \sum_{m=1}^{n}\mathbb P(A_m).
    }
\]

By assumption,
\[
    \sum_{m=1}^{\infty}\mathbb P(A_m)=\infty,
\]
and hence
\[
    \mathbb E(S_n)\longrightarrow\infty.
\]

We now use Chebyshev's inequality. For every $\delta>0$,
\[
\begin{aligned}
    \mathbb P\left(
        |S_n-\mathbb E S_n|
        >
        \delta\mathbb E S_n
    \right)
    &\leq
    \frac{\operatorname{Var}(S_n)}
         {\delta^2(\mathbb E S_n)^2}
    \\
    &\leq
    \frac{\mathbb E S_n}
         {\delta^2(\mathbb E S_n)^2}
    \\
    &=
    \frac{1}
         {\delta^2\mathbb E S_n}.
\end{aligned}
\]

Thus
\[
    \boxed{
    \mathbb P\left(
        |S_n-\mathbb E S_n|
        >
        \delta\mathbb E S_n
    \right)
    \leq
    \frac{1}{\delta^2\mathbb E S_n}
    \longrightarrow0.
    }
\]

Therefore,
\[
    \frac{S_n}{\mathbb E S_n}
    \xrightarrow{P}1.
\]

That is, we have convergence in probability.

To obtain almost sure convergence, we use a suitable subsequence.

\medskip

\noindent
\textbf{Construction of the subsequence.}

Since
\[
    \mathbb E S_n\longrightarrow\infty,
\]
for each $k\geq1$ define
\[
    n_k
    =
    \inf\left\{
        n:\mathbb E S_n\geq k^2
    \right\}.
\]

This is well-defined because $\mathbb E S_n\to\infty$.

By the definition of $n_k$,
\[
    \mathbb E S_{n_k}\geq k^2.
\]

Furthermore, since $n_k$ is the first index for which
$\mathbb E S_n\geq k^2$, we have
\[
    \mathbb E S_{n_k-1}<k^2.
\]

Since
\[
    \mathbb E X_m=\mathbb P(A_m)\leq1,
\]
we have
\[
    \mathbb E S_{n_k}
    =
    \mathbb E S_{n_k-1}
    +
    \mathbb E X_{n_k}
    <
    k^2+1.
\]

Consequently,
\[
    \boxed{
    k^2
    \leq
    \mathbb E S_{n_k}
    <
    k^2+1.
    }
\]

For convenience, define
\[
    T_k=S_{n_k}.
\]

Applying the previous Chebyshev estimate with $n=n_k$ gives
\[
\begin{aligned}
    \mathbb P\left(
        |T_k-\mathbb E T_k|
        >
        \delta\mathbb E T_k
    \right)
    &\leq
    \frac{1}
         {\delta^2\mathbb E T_k}.
\end{aligned}
\]

Since
\[
    \mathbb E T_k=\mathbb E S_{n_k}\geq k^2,
\]
we obtain
\[
    \boxed{
    \mathbb P\left(
        |T_k-\mathbb E T_k|
        >
        \delta\mathbb E T_k
    \right)
    \leq
    \frac{1}{\delta^2k^2}.
    }
\]

Therefore,
\[
\begin{aligned}
    \sum_{k=1}^{\infty}
    \mathbb P\left(
        |T_k-\mathbb E T_k|
        >
        \delta\mathbb E T_k
    \right)
    &\leq
    \frac{1}{\delta^2}
    \sum_{k=1}^{\infty}\frac1{k^2}
    \\
    &<\infty.
\end{aligned}
\]

By the first Borel--Cantelli lemma,
\[
    \mathbb P\left(
        |T_k-\mathbb E T_k|
        >
        \delta\mathbb E T_k
        \text{ i.o.}
    \right)=0.
\]

Thus, almost surely,
\[
    |T_k-\mathbb E T_k|
    \leq
    \delta\mathbb E T_k
\]
for all sufficiently large $k$.

Equivalently,
\[
    1-\delta
    \leq
    \frac{T_k}{\mathbb E T_k}
    \leq
    1+\delta
\]
eventually almost surely.

Since $\delta>0$ is arbitrary, we conclude that
\[
    \boxed{
    \frac{T_k}{\mathbb E T_k}
    \xrightarrow{\mathrm{a.s.}}1.
    }
\]

\medskip

\noindent
\textbf{Interpolation between the subsequence points.}

It remains to prove that
\[
    \frac{S_n}{\mathbb E S_n}
    \xrightarrow{\mathrm{a.s.}}1
\]
for the full sequence $n$, not only for the subsequence $n_k$.

Since
\[
    \mathbb E S_n
    =
    \sum_{m=1}^{n}\mathbb P(A_m)
\]
is nondecreasing in $n$, and $S_n$ is also nondecreasing because
$X_m=\mathbf{1}_{A_m}\geq0$, we have, whenever
\[
    n_k\leq n<n_{k+1},
\]
that
\[
    T_k
    =
    S_{n_k}
    \leq
    S_n
    \leq
    S_{n_{k+1}}
    =
    T_{k+1}.
\]

Also,
\[
    \mathbb E T_k
    =
    \mathbb E S_{n_k}
    \leq
    \mathbb E S_n
    \leq
    \mathbb E S_{n_{k+1}}
    =
    \mathbb E T_{k+1}.
\]

Therefore,
\[
    \frac{T_k}{\mathbb E T_{k+1}}
    \leq
    \frac{S_n}{\mathbb E S_n}
    \leq
    \frac{T_{k+1}}{\mathbb E T_k}.
\]

We rewrite the two bounds as
\[
\begin{aligned}
    \frac{T_k}{\mathbb E T_{k+1}}
    &=
    \frac{\mathbb E T_k}
         {\mathbb E T_{k+1}}
    \frac{T_k}{\mathbb E T_k},
    \\[6pt]
    \frac{T_{k+1}}{\mathbb E T_k}
    &=
    \frac{T_{k+1}}
         {\mathbb E T_{k+1}}
    \frac{\mathbb E T_{k+1}}
         {\mathbb E T_k}.
\end{aligned}
\]

Thus
\[
    \boxed{
    \frac{\mathbb E T_k}
         {\mathbb E T_{k+1}}
    \frac{T_k}{\mathbb E T_k}
    \leq
    \frac{S_n}{\mathbb E S_n}
    \leq
    \frac{T_{k+1}}
         {\mathbb E T_{k+1}}
    \frac{\mathbb E T_{k+1}}
         {\mathbb E T_k}.
    }
\]

We already know that
\[
    \frac{T_k}{\mathbb E T_k}
    \xrightarrow{\mathrm{a.s.}}1
\]
and
\[
    \frac{T_{k+1}}{\mathbb E T_{k+1}}
    \xrightarrow{\mathrm{a.s.}}1.
\]

It therefore remains only to show that
\[
    \frac{\mathbb E T_{k+1}}
         {\mathbb E T_k}
    \longrightarrow1.
\]

From the construction of $n_k$,
\[
    k^2
    \leq
    \mathbb E T_k
    <
    k^2+1
\]
and
\[
    (k+1)^2
    \leq
    \mathbb E T_{k+1}
    <
    (k+1)^2+1.
\]

Hence
\[
    \frac{\mathbb E T_{k+1}}
         {\mathbb E T_k}
    \leq
    \frac{(k+1)^2+1}{k^2}.
\]

But
\[
    \frac{(k+1)^2+1}{k^2}
    =
    \frac{k^2+2k+2}{k^2}
    =
    1+\frac{2}{k}+\frac{2}{k^2}
    \longrightarrow1.
\]

Also,
\[
    \frac{\mathbb E T_{k+1}}
         {\mathbb E T_k}
    \geq1,
\]
because $\mathbb E T_k$ is nondecreasing.

Therefore,
\[
    \boxed{
    \frac{\mathbb E T_{k+1}}
         {\mathbb E T_k}
    \longrightarrow1.
    }
\]

Consequently,
\[
    \frac{\mathbb E T_k}
         {\mathbb E T_{k+1}}
    \longrightarrow1.
\]

Hence, almost surely,
\[
    \frac{T_k}{\mathbb E T_{k+1}}
    \longrightarrow1
\]
and
\[
    \frac{T_{k+1}}{\mathbb E T_k}
    \longrightarrow1.
\]

Therefore, by the squeeze theorem,
\[
    \boxed{
    \frac{S_n}{\mathbb E S_n}
    \xrightarrow{\mathrm{a.s.}}1.
    }
\]

Finally, since
\[
    S_n=\sum_{m=1}^{n}\mathbf{1}_{A_m}
\]
and
\[
    \mathbb E S_n
    =
    \sum_{m=1}^{n}\mathbb P(A_m),
\]
we obtain
\[
    \boxed{
    \frac{\displaystyle\sum_{m=1}^{n}\mathbf{1}_{A_m}}
         {\displaystyle\sum_{m=1}^{n}\mathbb P(A_m)}
    \xrightarrow{\mathrm{a.s.}}1.
    }
\]

\end{proof}


% ============================================================
\lectureend
% ============================================================
